\section{System}
\paragraph{Acoustic front end.} The wearer is modelled as a binaural earbud pair (left/right
microphones 18\,cm apart on the interaural axis, free field, no head shadow; 16\,kHz). A 2-mic array
observes only the lateral angle, so front/back are confounded. Head rotation is rendered by
cross-fading image-source RIRs computed on a 5$^\circ$ yaw grid.
\paragraph{Pipeline.} Every 100\,ms the selection policy $\tau$ outputs the target's head-relative
azimuth $\theta$ (or ``unknown'') and which enrolled voice to condition on. Every hop, the streaming
extractor $f$ consumes $H$ new samples per microphone and emits $H$ output samples; the optional
personalisation $\rho$ multiplies the extractor's output spectrum by audiogram-dependent band gains in
the same STFT, adding no algorithmic latency; device volume and a hard safety limiter close every
playback path (Python, C++, Android). The C++ engine (ONNX Runtime with pre-bound I/O, no allocation
or locks on the audio callback) and the Android app (Oboe full-duplex, AAudio low-latency path) run the
same exported ONNX graph.
\paragraph{Latency accounting.} Algorithmic latency $= 2H +$ look-ahead (synthesis window support of
the asymmetric STFT); compute latency $=$ measured per-hop inference time (we use the p99); system
latency $=$ microphone-to-ear including the OS audio path and, for Bluetooth, the radio link. The
20\,ms constraint is checked on algorithmic $+$ p99 compute. System latency is \emph{not} measured
(no phone) and is reported only as labelled estimates from platform requirements and published claims
\citep{androidcdd15}.
